\documentclass[12pt]{article}
\usepackage[margin=1in]{geometry}
\usepackage{amsmath,amssymb}
\usepackage{graphicx}
\usepackage{booktabs}
\usepackage{hyperref}
\usepackage{natbib}
\usepackage{lineno}
\usepackage{setspace}

\doublespacing

\title{An Atlas of Brain--Bone Sympathetic Neural Circuits Mapped by Pseudorabies Viral Tracing}

\author{Author One$^{1}$, Author Two$^{1}$, Author Three$^{2}$ \\
\small $^{1}$Department of Neuroscience, University of Example, City, Country \\
\small $^{2}$Department of Orthopedic Surgery, University of Example, City, Country}

\date{}

\begin{document}

\maketitle

\begin{abstract}
The sympathetic nervous system (SNS) regulates bone metabolism, yet the brain regions directing this outflow remained uncharted. Here we provide the first comprehensive atlas of central SNS circuits projecting to bone using retrograde pseudorabies virus (PRV152) tracing from the mouse femur. Microinjection of EGFP-expressing PRV152 into five loci across the right femur periosteum and metaphysis of adult male mice (N=6; N=4 included in analysis) labeled neurons across 87 brain nuclei, sub-nuclei, and regions spanning six brain divisions: midbrain and pons, hypothalamus, medulla, forebrain, cerebral cortex, and thalamus. The raphe magnus (RMg) and periaqueductal gray (PAG) showed the highest infection levels. Bilateral brain labeling was observed from unilateral injections, and PRV152-infected neurons were confirmed in the intermediolateral cell column (IML) at T13--L2 spinal cord levels, validating the bone--SNS ganglia--IML--brain route. Negative controls (virus placed on bone surface) showed no EGFP signal. Comparison with prior rat bone marrow tracing data revealed overlapping but distinct circuits, with shared nodes including PAG, somatosensory cortex, MPOM, periventricular thalamic nucleus, PVH, lateral hypothalamic nucleus, and RPa. These findings establish the neuroanatomical foundation for understanding central regulation of bone physiology and bone pain.
\end{abstract}

\section{Introduction}

Bone is a richly innervated tissue. Histological studies have demonstrated abundant sympathetic nervous system (SNS) fibers in the periosteum and bone marrow, expressing markers including tyrosine hydroxylase (TH), dopamine beta-hydroxylase (DBH), neuropeptide Y (NPY), vesicular acetylcholine transporter (VAChT), and vasoactive intestinal polypeptide (VIP) \citep{bjurholm1991, mach2002, elefteriou2005}. These nerve fibers play critical roles in regulating bone formation and resorption through the release of norepinephrine and other neurotransmitters \citep{takeda2002, elefteriou2005}.

The central neural pathways that drive sympathetic outflow to bone have been the subject of growing interest. Leptin's anti-osteogenic actions are known to be mediated through the ventromedial hypothalamus (VMH) via the SNS \citep{ducy2000, takeda2002}, and recent work has identified PDE5A-containing neurons in specific brain nuclei linked to bone regulation \citep{huang2020}. However, a comprehensive mapping of all brain regions that participate in the central sympathetic outflow to bone has been lacking. Only one prior study attempted hierarchical circuit tracing from bone, using pseudorabies virus injection into rat femoral bone marrow, and that study covered a limited set of brain sites \citep{zhang2020}.

Pseudorabies virus (PRV) is a powerful tool for mapping polysynaptic neural circuits. Attenuated strains of PRV, such as the EGFP-expressing PRV152 (Bartha strain), travel exclusively in the retrograde direction from the site of peripheral injection, sequentially infecting neurons along the entire periphery-to-brain neuroaxis \citep{card2011, enquist2002}. This property makes PRV152 ideal for identifying all brain regions that contribute to the sympathetic innervation of a specific peripheral target.

Understanding the central circuits that direct sympathetic outflow to bone has broad implications. Bone pain, a major clinical problem associated with fractures, metastatic cancer, and osteoarthritis, is thought to involve both sensory afferents and sympathetic efferents \citep{mantyh2014}. The paraventricular hypothalamic nucleus (PVH)--periaqueductal gray (PAG)--raphe pallidus/raphe magnus (RPa/RMg) pathway has been implicated in descending pain modulation through regulation of enkephalin release in the spinal cord dorsal horn \citep{fields2004}. Identifying the full complement of brain regions in the bone--directed SNS circuit is essential for understanding the neural basis of bone pain and for developing targeted therapeutic interventions.

Here we report the results of PRV152 retrograde tracing from the mouse femur, providing the first comprehensive atlas of 87 brain nuclei, sub-nuclei, and regions across six brain divisions that send sympathetic efferent projections to bone.

\section{Methods}

\subsection{Animals and Housing}

Adult male mice (3--4 months old, N=6) were housed under a 12:12 hour light:dark cycle at $22 \pm 2$\textdegree C with ad libitum access to food and water. All procedures were approved by the institutional animal care and use committee.

\subsection{Surgical Procedure and Viral Injection}

Mice were anesthetized with isoflurane (2--3\% in oxygen). PRV152, an EGFP-expressing pseudorabies virus (Bartha strain; titer $4.7 \times 10^9$ pfu/mL), was microinjected into the right femur--tibia joint. Five injection loci were distributed across the bone metaphysis and periosteum, regions known to be enriched with SNS innervation (Table~\ref{tab:protocol}). Each locus received 150~nL of viral suspension, with a 60-second hold after each injection to minimize leakage. Animals were sacrificed 6 days post-injection based on pilot timing studies that determined the optimal survival time for labeling the complete neuroaxis without excessive infection.

\begin{table}[ht]
\centering
\caption{Detailed experimental protocol parameters for PRV152 tracing.}
\label{tab:protocol}
\begin{tabular}{ll}
\toprule
\textbf{Parameter} & \textbf{Value} \\
\midrule
Animals injected / included & 6 / 4 \\
Age & 3--4 months \\
Injection volume per locus & 150 nL \\
Number of injection loci & 5 \\
Viral titer & $4.7 \times 10^9$ pfu/mL \\
Post-injection hold time & 60 seconds \\
Survival time & 6 days \\
Section thickness & 25 $\mu$m \\
Primary antibody & Chicken anti-GFP (1:500) \\
Secondary antibody & Alexa Fluor 488 goat anti-chicken (1:200) \\
Blocking serum & 10\% normal goat serum \\
\bottomrule
\end{tabular}
\end{table}

\subsection{Tissue Processing and Immunofluorescence}

Animals were transcardially perfused with 4\% paraformaldehyde. Brains were post-fixed for 3--4 hours at 4\textdegree C, cryoprotected in 30\% sucrose in 0.1~M PBS, and sectioned at 25~$\mu$m on a freezing stage sliding microtome. Free-floating sections were blocked in 10\% normal goat serum, incubated overnight at 4\textdegree C with chicken anti-GFP primary antibody (1:500), and then incubated for 2 hours at room temperature with Alexa Fluor 488-conjugated goat anti-chicken secondary antibody (1:200). Sections were examined using a Nikon Eclipse Ni-E fluorescence microscope, and brain regions were identified using the Allen Mouse Brain Atlas.

\subsection{Animal Inclusion and Exclusion}

Four of six mice were included in the final analysis (Table~\ref{tab:inclusion}). Two mice were excluded due to over-infection, evidenced by widespread cloudy plaques across multiple brain regions, which prevents accurate quantification of region-specific labeling. The four included mice showed even infection across the neuroaxis without evidence of non-specific spread.

\begin{table}[ht]
\centering
\caption{Animal inclusion summary showing infection quality assessment.}
\label{tab:inclusion}
\begin{tabular}{lccc}
\toprule
\textbf{Mouse} & \textbf{Infection Quality} & \textbf{Included} \\
\midrule
1--4 & Even infection across neuroaxis & Yes \\
5 & Over-infected (widespread cloudy plaques) & No \\
6 & Over-infected (widespread cloudy plaques) & No \\
\bottomrule
\end{tabular}
\end{table}

\subsection{Negative Control}

A negative control experiment was performed in which PRV152 was placed on the bone surface rather than injected into the periosteum or metaphysis. Brain sections from negative control animals were examined for EGFP signal in PVH and RPa, two regions known to be labeled in successful injections.

\section{Results}

\subsection{Validation of PRV152 Tracing}

The negative control (PRV152 placed on bone surface) produced no detectable EGFP signal in either PVH or RPa, confirming that viral injection into the periosteum or metaphysis is required for retrograde labeling. Experimental injections produced robust EGFP immunoreactivity in both PVH and RPa, as well as PRV152-infected neurons in the intermediolateral cell column (IML) at T13--L2 spinal cord levels. IML labeling validates the expected route of retrograde infection: bone $\rightarrow$ SNS ganglia $\rightarrow$ IML (spinal cord) $\rightarrow$ brain.

\subsection{Comprehensive Atlas of Brain Regions}

PRV152-labeled neurons were identified across 87 brain nuclei, sub-nuclei, and regions spanning six brain divisions (Table~\ref{tab:divisions}). The midbrain and pons, hypothalamus, and medulla contained the highest number of labeled regions, consistent with these divisions housing key autonomic and pre-autonomic neuronal populations. The forebrain, cerebral cortex, and thalamus also contained labeled regions, indicating that higher-order brain areas participate in the sympathetic outflow to bone.

\begin{table}[ht]
\centering
\caption{Distribution of PRV152-labeled brain regions across six brain divisions, with key nuclei showing highest infection levels.}
\label{tab:divisions}
\begin{tabular}{llll}
\toprule
\textbf{Brain Division} & \textbf{Key Nuclei} & \textbf{Infection Level} \\
\midrule
Midbrain \& Pons & PAG & High \\
Hypothalamus & LH, PVH & High / Moderate-High \\
Medulla & RMg, RPa & High \\
Forebrain & MPOM & High \\
Cerebral Cortex & S1HL (somatosensory) & High \\
Thalamus & Periventricular nucleus & High \\
\midrule
\multicolumn{2}{l}{\textbf{Total nuclei/sub-nuclei/regions}} & \textbf{87} \\
\bottomrule
\end{tabular}
\end{table}

\subsection{Site-Specific Variation in Infection Levels}

Significant variation in PRV152 labeling intensity was observed across brain regions. The raphe magnus (RMg) of the medulla and the periaqueductal gray (PAG) of the midbrain showed the highest levels of EGFP-labeled neurons. Other regions with high infection levels included the lateral hypothalamus (LH), raphe pallidus (RPa), medial preoptic nucleus medial part (MPOM), primary somatosensory cortex hindlimb area (S1HL), periventricular thalamic nucleus, and paraventricular hypothalamic nucleus (PVH). The high labeling in RMg and PAG is consistent with the known roles of these regions in autonomic regulation and descending pain modulation.

\subsection{Bilateral Labeling Pattern}

Despite unilateral injection into the right femur, PRV152-labeled neurons were observed bilaterally in the brain with no noticeable hemispheric dominance. This bilateral pattern is consistent with prior PRV tracing studies of SNS innervation of peripheral tissues, including fat depots, and likely reflects the bilateral organization of sympathetic pre-ganglionic neurons and their supraspinal inputs.

\subsection{Comparison with Prior Bone Marrow Tracing}

Comparison of our femur tracing results with a prior study that traced SNS circuits from rat femoral bone marrow revealed overlapping but distinct circuit compositions (Table~\ref{tab:comparison}). Seven key brain regions were labeled in both studies: PAG, somatosensory cortex, MPOM, periventricular thalamic nucleus, PVH, lateral hypothalamic nucleus, and RPa. However, additional regions were uniquely identified in each study, suggesting that separate site-specific SNS circuits may project to the femur versus the femoral bone marrow.

\begin{table}[ht]
\centering
\caption{Comparison of shared brain regions between femur (this study, mouse) and bone marrow (prior study, rat) PRV tracing.}
\label{tab:comparison}
\begin{tabular}{lccc}
\toprule
\textbf{Brain Region} & \textbf{Femur (this study)} & \textbf{Bone Marrow (prior)} & \textbf{Shared} \\
\midrule
PAG (midbrain) & Yes & Yes & Yes \\
Somatosensory cortex & Yes & Yes & Yes \\
MPOM (forebrain) & Yes & Yes & Yes \\
Periventricular nucleus & Yes & Yes & Yes \\
PVH (hypothalamus) & Yes & Yes & Yes \\
Lateral hypothalamic nucleus & Yes & Yes & Yes \\
RPa (medulla) & Yes & Yes & Yes \\
\bottomrule
\end{tabular}
\end{table}

\subsection{Pain Circuit Relevance}

The brain regions identified in this atlas include key nodes of the descending pain modulation pathway. The PVH--PAG--RPa/RMg axis represents a well-characterized circuit in which PVH pre-autonomic neurons project to PAG, which in turn projects to RPa and RMg in the medulla. These medullary nuclei project to the dorsal horn of the spinal cord, where they regulate enkephalin release and thereby modulate pain perception. The labeling of all nodes of this pathway in our bone tracing study provides a neuroanatomical basis for understanding how the brain may modulate bone pain through descending sympathetic circuits.

\section{Discussion}

This study provides the first comprehensive atlas of brain regions sending sympathetic nervous system efferent projections to bone, identifying 87 nuclei, sub-nuclei, and regions across all six major brain divisions. The breadth of this central outflow network is striking, encompassing not only brainstem autonomic centers that would be expected to participate in SNS regulation, but also higher-order regions in the cerebral cortex, thalamus, and forebrain. This widespread distribution suggests that bone sympathetic tone is subject to diverse modulatory inputs spanning sensory, cognitive, and emotional domains.

The particularly high PRV152 labeling in the periaqueductal gray and raphe magnus is notable for several reasons. Both regions are established components of the descending pain modulation system \citep{fields2004, basbaum2009}, and their strong labeling suggests robust polysynaptic connectivity with bone-innervating sympathetic neurons. The PAG is also a critical node for coordinating autonomic, somatomotor, and behavioral responses to stressors \citep{bandler2000}. The convergence of pain modulation and autonomic regulation in these circuits may explain the bidirectional relationship between bone pain and sympathetic nervous system activity observed clinically.

The bilateral brain labeling from unilateral femur injection is consistent with the known bilateral organization of sympathetic circuits and has been observed in PRV tracing studies of other peripheral targets \citep{stanley2010}. This pattern implies that bone sympathetic regulation is coordinated across both hemispheres, rather than being strictly lateralized to the ipsilateral side.

The identification of somatosensory cortex (S1HL) among the highly labeled regions is intriguing. While the somatosensory cortex is primarily associated with processing sensory afferent information, its presence in a retrograde SNS tracing study suggests the existence of cortico-autonomic pathways that may provide top-down modulatory influence on bone sympathetic tone. Such pathways could mediate the effects of somatosensory stimulation on bone metabolism, a relationship that has been observed experimentally but whose neural substrate was previously unknown \citep{sample2015}.

The comparison with prior rat bone marrow tracing data reveals both shared and distinct circuit elements. The seven shared nodes---PAG, somatosensory cortex, MPOM, periventricular nucleus, PVH, lateral hypothalamic nucleus, and RPa---likely represent core components of the central sympathetic outflow to bone that are conserved across bone compartments and species. The differences between the two studies may reflect genuine site-specific innervation patterns, species differences between mouse and rat, or methodological variations.

Several limitations should be acknowledged. The sample size of four mice included in the analysis, while typical for viral tracing studies, limits our ability to quantify the reliability of labeling in each region. The exclusion of two mice due to over-infection underscores the challenge of controlling viral spread in transneuronal tracing experiments. The 6-day survival time was optimized for labeling the complete neuroaxis, but different survival times might reveal additional or fewer labeled regions at different levels of the circuit hierarchy. Finally, PRV152 tracing does not distinguish between direct and indirect polysynaptic pathways, so the number of synaptic relays between the brain and the bone cannot be determined from these data alone.

\section{Conclusion}

We have established a comprehensive atlas of 87 brain regions that send sympathetic nervous system efferent projections to bone in the mouse. The atlas spans six brain divisions, with the raphe magnus and periaqueductal gray showing the highest levels of retrograde viral labeling. These findings provide the neuroanatomical foundation for understanding central regulation of bone physiology, including bone metabolism, bone pain, and the interface between the nervous and skeletal systems.

\bibliographystyle{plainnat}
\bibliography{references}

\end{document}
